\documentclass[11pt]{article}
\usepackage{amsmath, amssymb, amsthm, graphicx}

\title{Kaprekar Dynamics on Digit-Bag Quotients:\\A Finite Spectral Graph Reduction}
\author{}
\date{}

\newtheorem{theorem}{Theorem}
\newtheorem{definition}{Definition}
\newtheorem{proposition}{Proposition}

\begin{document}

\maketitle

\begin{abstract}
We study the Kaprekar transformation on $d$-digit integers as a finite dynamical system. 
We introduce a permutation-invariant quotient by digit multisets (digit-bags), reducing the full state space of size $10^d$ to a structured quotient graph. 
For $d=5$, we compute the exact quotient system of 2,002 nodes and show that it forms a deterministic functional graph with four attractor cycles and pure basin decomposition. 
We further analyze the induced Laplacian structure, spectral properties, and effective resistance geometry of the quotient system.
\end{abstract}

\section{Introduction}

The Kaprekar map is defined on finite digit strings by sorting digits in descending and ascending order and subtracting:
\[
T(n) = \text{desc}(n) - \text{asc}(n).
\]

This induces a deterministic functional graph on $\{0,\dots,10^d-1\}$.

While classical results focus on fixed points (e.g. 6174), higher-dimensional cases exhibit richer basin and cycle structures.

We study the system through a permutation-invariant quotient space.

\section{System Definition}

\begin{definition}[Kaprekar Map]
Let $X_d = \{0,\dots,10^d-1\}$. Define:
\[
T(x) = \text{desc}(x) - \text{asc}(x).
\]
\end{definition}

\begin{definition}[Digit-Bag Quotient]
Define an equivalence relation:
\[
x \sim y \iff \text{multiset}(x) = \text{multiset}(y).
\]
Let the quotient space be:
\[
B_d = X_d / \sim.
\]
\end{definition}

\section{Induced Quotient Dynamics}

\begin{proposition}
The Kaprekar map induces a well-defined map $\tilde{T}: B_d \to B_d$.
\end{proposition}

\begin{proof}
Sorting and subtraction commute with digit permutation; hence if $x \sim y$, then $T(x) \sim T(y)$.
\end{proof}

Thus the system reduces to a deterministic functional graph on $|B_d|$ nodes.

For $d=5$, we have $|B_5| = 2002$.

\section{Graph Structure}

The quotient system defines a directed graph $G_B = (B_d, E)$ with:
\[
E = \{(b, \tilde{T}(b))\}.
\]

\begin{theorem}[Functional Graph Decomposition]
Every finite functional graph decomposes into disjoint cycles with in-trees feeding into them.
\end{theorem}

For $d=5$, the system contains:
\begin{itemize}
\item 4 attractor cycles
\item maximum depth = 5
\item average depth ≈ 2.08
\end{itemize}

\section{Basin Structure}

\begin{definition}
A basin is the set of nodes flowing into a given cycle.
\end{definition}

\begin{proposition}
For $d=5$, every digit-bag lies entirely within a single basin.
\end{proposition}

This implies a perfect partition:
\[
B_d = \bigsqcup_{k=1}^4 \mathcal{B}_k.
\]

No digit-bag intersects multiple basins.

\section{Cycle Structure (d = 5)}

Empirically computed cycles:

\begin{itemize}
\item Period-1 cycle (fixed point)
\item One period-2 cycle
\item Two period-4 cycles
\end{itemize}

Basin sizes:

\[
(10,\; 48320,\; 3190,\; 48480).
\]

\section{Spectral Structure}

Define transition matrix:
\[
P_{ij} =
\begin{cases}
1 & j = T(i) \\
0 & \text{otherwise}
\end{cases}
\]

Laplacian:
\[
L = I - P.
\]

\begin{proposition}
The spectrum of $P$ is determined by cycle structure and nilpotent tree components.
\end{proposition}

Eigenvalues on cycles correspond to roots of unity.

Tree structure contributes nilpotent Jordan blocks.

\section{Effective Resistance Geometry}

Define:
\[
R_{ij} = (e_i - e_j)^T L^+ (e_i - e_j).
\]

This induces a geometric embedding of the quotient graph.

Empirically:
\begin{itemize}
\item resistance correlates weakly with depth $\tau$
\item bottlenecks occur near basin boundaries
\end{itemize}

\section{Symmetry Analysis}

Define digit-wise complement:
\[
d \mapsto 9 - d.
\]

\begin{proposition}
The 9-complement preserves basin structure on $B_d$.
\end{proposition}

However:
\begin{itemize}
\item no basin-swapping involution exists
\item no self-complementary bags exist
\end{itemize}

Thus symmetry acts trivially on basin partition.

\section{Main Structural Insight}

\begin{theorem}[Quotient Dynamics Principle]
The Kaprekar map on $X_d$ factors through the digit-bag quotient $B_d$, producing a deterministic finite dynamical system whose full dynamical structure is encoded in a 2002-node graph for $d=5$.
\end{theorem}

\section{Conclusion}

The Kaprekar system is not chaotic in a classical sense but is a permutation-equivariant finite dynamical system. Its true structure emerges only under quotienting by digit permutations, revealing a small deterministic graph with rich basin geometry and spectral structure.

\section{Future Work}

\begin{itemize}
\item full Laplacian spectrum of quotient graph
\item comparison with lifted full-state spectrum
\item multiscale coarsening of basin structure
\item resistance geometry classification of transitions
\end{itemize}

\end{document}
